\documentclass[../DoAn.tex]{subfiles}
\begin{document}

\section{Thiết kế kiến trúc}
\label{section:4.1}

\subsection{Lựa chọn kiến trúc phần mềm}
\label{subsection:4.1.1}

Hệ thống gồm năm nhóm thành phần chính: kênh truy cập, backend Spring Boot, lưu trữ (PostgreSQL và file system), chatbot runtime FastAPI và dịch vụ LLM bên ngoài. Giao diện web dùng cho quản trị hệ thống, quản trị cửa hàng và thao tác chat. Các kênh Messenger và Telegram gửi tin nhắn về backend thông qua webhook — phần hạ tầng webhook thuộc về nền tảng bên ngoài, backend chỉ cấu hình endpoint để nhận dữ liệu.

Backend Spring Boot tiếp nhận yêu cầu từ giao diện web qua HTTP API và từ Messenger/Telegram qua endpoint webhook. MessengerWebhookController và TelegramWebhookController nhận POST tại \texttt{/webhook/messenger} và \texttt{/webhook/telegram/\{secretPath\}}. Backend xử lý xác thực, phân quyền, quản lý nguồn dữ liệu sản phẩm, quản lý phiên bản knowledge base, lưu hội thoại, lưu khách hàng tiềm năng và lưu yêu cầu mua hàng.

Chatbot runtime FastAPI xử lý phần truy xuất tri thức và sinh phản hồi. Khi nhận yêu cầu từ backend, runtime đọc trực tiếp file knowledge base từ thư mục trên file system (theo đường dẫn \texttt{kb\_dir} được backend cung cấp), truy xuất các nội dung liên quan đến câu hỏi, tạo prompt và gửi yêu cầu đến Claude API (Anthropic) để sinh phản hồi. Trường hợp API không khả dụng hoặc retrieval không đủ context, hệ thống dùng luật dự phòng bằng template-based response. Kết quả trả về cho backend là nội dung phản hồi và các tín hiệu phục vụ xử lý hội thoại.

PostgreSQL lưu dữ liệu vận hành của hệ thống, gồm tenant, hội thoại, message, product\_dataset, artifact metadata, lead và purchase\_request. Các bảng này đều có trường \texttt{tenant\_id} để phân tách theo cửa hàng. KB Artifact được lưu trên file system dưới dạng thư mục chứa file catalog, rag products và manifest; backend ghi các file này khi build và lưu đường dẫn (artifactPath, kb\_dir) xuống database.

Giao diện quản trị và giao diện cửa hàng được xây dựng bằng HTML, CSS và JavaScript thuần. Các giao diện gọi API backend để quản lý cửa hàng, nguồn tri thức, phiên bản knowledge base, hội thoại, khách hàng tiềm năng, yêu cầu mua hàng và trạng thái của chatbot.

\subsection{Thiết kế tổng quan}
\label{subsection:4.1.2}

Hình~\ref{fig:ch4_overall_multitenant_architecture} mô tả các thành phần chính của hệ thống và hướng trao đổi dữ liệu giữa các thành phần theo kiến trúc phân tầng dọc. Tầng trên cùng là các kênh truy cập: Web Chat và Admin UI chạy trên trình duyệt, Messenger và Telegram là các nền tảng gửi tin nhắn qua webhook — phần hạ tầng webhook này thuộc về nền tảng bên ngoài, không phải thành phần do hệ thống xây dựng. Tầng giữa là Spring Boot backend tiếp nhận từ cả ba kênh.

\begin{figure}[H]
    \centering
    \includegraphics[width=0.95\textwidth]{Hinhve/Hinh4_1_kien_truc_tong_the_multitenant_rag.png}
    \caption{Kiến trúc tổng thể multi-tenant RAG chatbot}
    \label{fig:ch4_overall_multitenant_architecture}
\end{figure}

Trong Hình~\ref{fig:ch4_overall_multitenant_architecture}, Spring Boot backend tiếp nhận yêu cầu từ các kênh truy cập. Với Web Chat và Admin UI, backend nhận HTTP API request. Với Messenger và Telegram, MessengerWebhookController và TelegramWebhookController nhận dữ liệu qua endpoint \texttt{/webhook/messenger} và \texttt{/webhook/telegram/\{secretPath\}}. Backend xác định cửa hàng đang thao tác, kiểm tra quyền truy cập, lưu hội thoại, quản lý nguồn dữ liệu sản phẩm, quản lý phiên bản knowledge base, quản lý khách hàng tiềm năng và yêu cầu mua hàng. Với các yêu cầu hội thoại, backend xác định cửa hàng, kênh gửi tin nhắn, phiên hội thoại và chế độ hội thoại trước khi chuyển yêu cầu sang chatbot runtime.

Chatbot runtime được tổ chức theo cửa hàng. Khi nhận yêu cầu từ backend, runtime đọc trực tiếp file knowledge base từ thư mục trên file system (theo đường dẫn \texttt{kb\_dir} được backend cung cấp thông qua TenantKbVersion), truy xuất các nội dung liên quan đến câu hỏi. Sau bước truy xuất, runtime tạo prompt và gửi đến Claude API (Anthropic). Phản hồi từ Claude API được runtime nhận về để xử lý. Trường hợp API không khả dụng hoặc không đủ context để trả lời, hệ thống dùng luật dự phòng dạng template-based do hệ thống tự xây dựng. Kết quả phản hồi được trả lại backend để lưu vào hội thoại và gửi về đúng kênh người dùng đã sử dụng.

PostgreSQL lưu dữ liệu vận hành của hệ thống, gồm dữ liệu cửa hàng, tài khoản, hội thoại, khách hàng, product\_dataset, artifact metadata, khách hàng tiềm năng và yêu cầu mua hàng. KB Artifact được lưu trên file system dưới dạng thư mục chứa catalog, rag products và manifest; backend ghi dữ liệu vào thư mục này khi build artifact và ghi đường dẫn (artifactPath, kb\_dir) vào database. Khi cần retrieval, chatbot runtime đọc trực tiếp file KB từ thư mục theo đường dẫn kb\_dir. Khi đọc hoặc ghi dữ liệu, backend luôn xác định dữ liệu thuộc cửa hàng nào trước khi thực hiện thao tác. Nhờ đó, các cửa hàng có thể dùng chung hệ thống nhưng dữ liệu hội thoại, khách hàng và nghiệp vụ bán hàng không bị dùng lẫn.

Knowledge base được lưu tách khỏi dữ liệu hội thoại và dữ liệu nghiệp vụ. Dữ liệu sản phẩm sau khi thu thập không được đưa trực tiếp vào chatbot, mà đi qua bước kiểm tra chất lượng, tạo artifact và kích hoạt phiên bản knowledge base cho cửa hàng. Khi một phiên bản knowledge base mới được kích hoạt, runtime của cửa hàng đó được làm mới để lần chat tiếp theo sử dụng đúng nguồn tri thức đang được chọn.

Cả ba kênh Web, Messenger và Telegram đều đi qua cùng một luồng xử lý trong chatbot runtime. Sau khi backend tiếp nhận tin nhắn, hệ thống xác định tenant, conversation và channel, sau đó chuyển sang runtime để retrieval và sinh phản hồi. Runtime không phân biệt kênh gửi tin nhắn là Web, Messenger hay Telegram — chỉ sử dụng tenant, conversation và active KB version để xử lý. Kết quả phản hồi được backend gửi về đúng kênh ban đầu.

\section{Thiết kế chi tiết}
\label{section:4.2}

\subsection{Thiết kế xử lý theo tenant}
\label{subsection:4.2.1}

Trong phạm vi hệ thống, tenant đại diện cho một cửa hàng sử dụng nền tảng. Mỗi tenant có tập dữ liệu riêng gồm chatbot, conversation, message, lead, purchase request, customer identity và active KB version. \texttt{tenant\_id} được dùng làm ranh giới kiểm soát xuyên suốt từ API, repository, KB version đến chatbot runtime. Khi backend nhận request, TenantContext xác định tenant hiện hành trước khi repository truy vấn dữ liệu — mọi câu lệnh truy vấn đều gắn điều kiện \texttt{WHERE tenant\_id = ?}.

Với platform admin, giao diện quản trị cho phép chọn tenant khi cần thao tác dữ liệu của một cửa hàng cụ thể. Backend kiểm tra quyền platform admin trước khi chấp nhận tenant được chọn. Với tenant admin và staff, backend xác định tenant dựa trên tài khoản, token hoặc tenant context hiện hành — tenant không được truyền tùy ý từ frontend. \texttt{SecurityConfig} kiểm tra vai trò người dùng, nếu là tenant admin thì chỉ cho phép truy cập dữ liệu thuộc tenant của họ.

Tenant được thiết kế độc lập với knowledge base. Một tenant có thể chưa có KB active; một KB Artifact có thể tồn tại độc lập trước khi được bind vào tenant. Khi tenant sử dụng một KB Artifact, backend tạo TenantKbBinding và TenantKbVersion tương ứng. Trên tenant, trường \texttt{activeKbVersionId} là nguồn chính để chatbot runtime xác định KB dir. Trường \texttt{tenant.kbDir} chỉ đóng vai trò fallback nếu tenant chưa có active KB version.

Khi bind KB mới, publish KB version mới hoặc unbind KB khỏi tenant, backend evict runtime của tenant đó. Runtime evict có nghĩa là runtime hiện tại bị xoá khỏi bộ nhớ, chat request tiếp theo sẽ lazy-load runtime mới dựa trên \texttt{activeKbVersionId} hiện hành. Nhiều conversation trong cùng tenant dùng chung runtime, nhưng runtime luôn theo đúng KB version đang active của tenant.

\subsection{Thiết kế Product Dataset, KB Artifact và Tenant Binding}
\label{subsection:4.2.2}

Pipeline offline được thiết kế để tách quá trình thu thập, kiểm tra, build và publish knowledge base thành các bước riêng, có thể truy vết trước khi tenant sử dụng. Dữ liệu crawl thô không được đưa trực tiếp vào chatbot runtime.

Hình~\ref{fig:ch4_offline_product_dataset_kb_artifact} mô tả luồng offline xây dựng Product Dataset và KB Artifact. Nguồn dữ liệu đầu vào có thể là URL sản phẩm, sitemap hoặc thư mục dữ liệu đã materialize. Sau khi crawl hoặc import, dữ liệu được đưa qua quality gate. Nếu đạt, dữ liệu được đăng ký thành ProductDataset trong database. Tiếp theo, backend chạy build để tạo ProductDatasetArtifact trên file system.

\begin{figure}[H]
    \centering
    \includegraphics[width=0.95\textwidth]{Hinhve/Hinh4_2_offline_product_dataset_kb_artifact.png}
    \caption{Luồng offline xây dựng Product Dataset và KB Artifact}
    \label{fig:ch4_offline_product_dataset_kb_artifact}
\end{figure}

Trong luồng ở Hình~\ref{fig:ch4_offline_product_dataset_kb_artifact}, ProductDataset là bản ghi trong database đại diện cho dataset gốc sau khi materialize. Dataset chứa đường dẫn tới các file trên file system: catalog, rag products, manifest, crawl report và quality audit. ProductDatasetArtifact là kết quả build KB từ ProductDataset — artifact chứa dữ liệu đã sẵn sàng phục vụ retrieval, nhưng chưa thuộc về tenant nào.

Quality gate được đặt trước bước register dataset. Các kiểm tra chính gồm: số lượng sản phẩm trong manifest có khớp với file thực tế hay không, số dòng catalog và rag có khớp hay không, có lỗi mojibake trong title/rag title không, tỷ lệ URL trùng lặp, độ phủ giá, độ phủ chất liệu, độ phủ kích thước và mức độ nghi ngờ category. Nếu quality gate phát hiện lỗi nghiêm trọng, dataset không được register; nếu lỗi nhỏ, dataset vẫn được register với trạng thái warn.

Sau khi ProductDatasetArtifact đạt trạng thái sẵn sàng, tenant có thể sử dụng artifact thông qua cơ chế binding. Hình~\ref{fig:ch4_kb_binding_unbinding} mô tả luồng bind/unbind.

\begin{figure}[H]
    \centering
    \includegraphics[width=0.95\textwidth]{Hinhve/Hinh4_3_bind_unbind_kb_tenant.png}
    \caption{Luồng bind/unbind KB Artifact với tenant}
    \label{fig:ch4_kb_binding_unbinding}
\end{figure}

Khi bind, backend tạo hoặc cập nhật TenantKbBinding giữa tenant và artifact. Sau đó backend tạo TenantKbVersion trỏ tới \texttt{artifact\_path}, cập nhật \texttt{activeKbVersionId} trên tenant và evict runtime. Request chat tiếp theo sẽ lazy-load runtime dựa trên active KB version mới.

Khi unbind, binding được đánh dấu không active, thời điểm unbind được lưu lại, \texttt{activeKbVersionId} trên tenant bị xoá và runtime tenant bị evict. Nếu tenant còn \texttt{kbDir} fallback, runtime có thể dùng đường dẫn này; nếu không có, tenant được xem như chưa có KB active.

Cơ chế binding hỗ trợ hai chính sách cập nhật. \texttt{AUTO\_USE\_LATEST}: khi dataset có artifact mới, các tenant đang active binding với dataset đó được cập nhật sang version mới và evict runtime. \texttt{MANUAL\_FIXED\_ARTIFACT}: tenant không tự chuyển sang artifact mới nếu chưa có thao tác quản trị.

\subsection{Thiết kế luồng RAG cho chế độ tenant\_sales}
\label{subsection:4.2.3}

Chế độ \texttt{tenant\_sales} là luồng chatbot bán hàng theo cửa hàng. Mục tiêu của luồng này là bảo đảm mỗi câu trả lời được sinh ra dựa trên đúng tenant, đúng conversation và đúng active KB version. Hình~\ref{fig:ch4_online_tenant_sales_rag} mô tả luồng xử lý online RAG cho chế độ này.

\begin{figure}[H]
    \centering
    \includegraphics[width=0.95\textwidth]{Hinhve/Hinh4_4_online_tenant_sales_rag.png}
    \caption{Luồng online RAG cho chế độ tenant\_sales}
    \label{fig:ch4_online_tenant_sales_rag}
\end{figure}

Trong Hình~\ref{fig:ch4_online_tenant_sales_rag}, luồng bắt đầu khi người dùng gửi tin nhắn từ Web Chat, Messenger hoặc Telegram. Backend tiếp nhận, xác định tenant, chatbot, conversation và channel tương ứng, đồng thời lưu tin nhắn vào bảng message. Tiếp theo, backend resolve active KB version: nếu tenant có \texttt{activeKbVersionId}, runtime sử dụng KB dir từ TenantKbVersion; nếu không, hệ thống xét đến fallback \texttt{tenant.kbDir}.

Sau khi xác định KB, backend gửi yêu cầu chat đến FastAPI runtime. Runtime thực hiện retrieval top-k từ KB của tenant (đọc file từ thư mục kb\_dir trên file system), xây dựng context, kết hợp với lịch sử hội thoại, mode rule và câu hỏi hiện tại để tạo prompt, sau đó gửi đến Claude API (Anthropic). Nếu API không phản hồi hoặc retrieval không đủ context, runtime dùng rule/template fallback do hệ thống tự xây dựng.

Kết quả từ Claude API hoặc fallback được backend lưu vào database dưới dạng assistant message. Nếu phản hồi chứa tín hiệu nghiệp vụ (ý định mua, thông tin liên hệ), backend kiểm tra các điều kiện: mode phải là \texttt{tenant\_sales}, conversation thuộc tenant hợp lệ, dữ liệu chưa trùng lặp. Chỉ khi tất cả điều kiện thỏa mãn, backend mới tạo lead hoặc purchase request. Cuối cùng, phản hồi được gửi về đúng kênh ban đầu của người dùng.

Hình~\ref{fig:ch4_runtime_lifecycle} mô tả vòng đời runtime sau khi publish KB version mới.

\begin{figure}[H]
    \centering
    \includegraphics[width=0.95\textwidth]{Hinhve/Hinh4_5_runtime_lifecycle_active_version.png}
    \caption{Vòng đời runtime chatbot theo tenant sau khi publish KB version}
    \label{fig:ch4_runtime_lifecycle}
\end{figure}

Trong Hình~\ref{fig:ch4_runtime_lifecycle}, publish KB version mới kích hoạt các bước sau. Bước 1: admin publish version. Bước 2: backend cập nhật \texttt{activeKbVersionId} trên tenant. Bước 3: backend evict runtime hiện tại — runtime cũ bị xoá khỏi bộ nhớ. Trạng thái runtime-status lúc này: \texttt{desired} trỏ tới version mới, \texttt{running} rỗng, \texttt{in\_sync=false}. Bước 4: runtime-status được sync với trạng thái hiện tại. Bước 5: \texttt{desired} được cập nhật trên runtime-status. Bước 6: khi có chat request tiếp theo, runtime được spawn lại dựa trên \texttt{activeKbVersionId} mới. Sau khi runtime chạy đúng version, runtime-status chuyển về \texttt{running} có giá trị, \texttt{process\_alive=true}, \texttt{in\_sync=true}.

\subsection{Thiết kế ba chế độ hội thoại}
\label{subsection:4.2.4}

Ba chế độ hội thoại được dùng để kiểm soát phạm vi dữ liệu, cách build prompt và tác động nghiệp vụ sau phản hồi. Hình~\ref{fig:ch4_chat_mode_control} mô tả ba mode: tenant\_sales, general\_compare và market\_price. Mỗi mode có nguồn dữ liệu riêng và tác động nghiệp vụ riêng.

\begin{figure}[H]
    \centering
    \includegraphics[width=0.95\textwidth]{Hinhve/Hinh4_6_kiem_soat_ba_mode_hoi_thoai.png}
    \caption{Kiểm soát ba chế độ hội thoại và tác động nghiệp vụ}
    \label{fig:ch4_chat_mode_control}
\end{figure}

Bảng~\ref{tab:ch4_chat_mode_comparison} so sánh ba mode. Tin nhắn từ người dùng được backend xác định mode dựa trên chatbot instance. Mỗi mode sử dụng nguồn dữ liệu khác nhau và có cơ chế kiểm soát tác động nghiệp vụ riêng.

\begin{table}[H]
\centering
\caption{So sánh ba chế độ hội thoại trong hệ thống}
\label{tab:ch4_chat_mode_comparison}
\resizebox{\textwidth}{!}{
\begin{tabular}{|p{3cm}|p{4cm}|p{4cm}|p{3.5cm}|p{4cm}|}
\hline
\textbf{Mode} & \textbf{Mục đích} & \textbf{Nguồn dữ liệu} & \textbf{Tác động nghiệp vụ} & \textbf{Cơ chế kiểm soát} \\
\hline
\texttt{tenant\_sales} & Tư vấn bán hàng cho một cửa hàng cụ thể. & Tenant-bound KB, active KB version, artifact đã bind với tenant. & Có thể tạo lead hoặc purchase request nếu đủ điều kiện. & Backend kiểm tra tenant, conversation, mode, contact, xác nhận mua và chống tạo trùng. \\
\hline
\texttt{general\_compare} & Tư vấn hoặc so sánh sản phẩm ở phạm vi chung. & Dữ liệu tham chiếu/general corpus được phép sử dụng. & Không tạo lead hoặc purchase request. & Chỉ phục vụ tư vấn/so sánh; nếu người dùng chỉ rõ source thì lọc theo source tương ứng. \\
\hline
\texttt{market\_price} & Tham khảo giá thị trường hoặc khoảng giá sản phẩm. & Dữ liệu giá hoặc observation được hệ thống quản lý. & Không tạo lead hoặc purchase request. & Không để LLM tự bịa khoảng giá; kết quả cần dựa trên dữ liệu quan sát hoặc provider có cấu trúc. \\
\hline
\end{tabular}
}
\end{table}

Trong phạm vi triển khai hiện tại, \texttt{tenant\_sales} là mode có đầy đủ implementation: tenant-bound KB, sales flow và purchase guard. \texttt{general\_compare} và \texttt{market\_price} phục vụ tư vấn chung hoặc tham chiếu giá — các luồng này chưa hoàn thiện general data layer và thống kê thị trường.

\subsection{Thiết kế nhận diện khách hàng liên kênh}
\label{subsection:4.2.5}

Hệ thống hỗ trợ ba kênh hội thoại: Web Chat, Messenger và Telegram. Mỗi kênh có external identity riêng: Web Chat dùng \texttt{userExternalId} hoặc \texttt{session\_id}; Messenger dùng \texttt{pageId} và \texttt{senderId}; Telegram dùng \texttt{chatId} hoặc \texttt{userId}. Hình~\ref{fig:ch4_cross_channel_identity} mô tả cơ chế nhận diện khách hàng qua các kênh này.

\begin{figure}[H]
    \centering
    \includegraphics[width=0.95\textwidth]{Hinhve/Hinh4_7_cross_channel_identity.png}
    \caption{Nhận diện khách hàng liên kênh Web, Messenger và Telegram}
    \label{fig:ch4_cross_channel_identity}
\end{figure}

Trong Hình~\ref{fig:ch4_cross_channel_identity}, ba kênh gửi định danh đến CustomerIdentityService. Service này tìm hoặc tạo bản ghi trong bảng \texttt{customer\_identities} theo tenant, channel và external id. Nếu người dùng cung cấp phone/email, service chuẩn hóa và liên kết các identity về cùng một \texttt{UnifiedCustomer} trong bảng \texttt{unified\_customers}. Conversation được gắn \texttt{unified\_customer\_id} để truy vết lịch sử qua nhiều kênh.

Hệ thống không merge khách hàng chỉ dựa trên \texttt{displayName}. Việc merge chỉ thực hiện khi có strong identifiers (phone/email đã chuẩn hóa). Tenant isolation được duy trì: cùng phone/email nhưng khác tenant không bị trộn vào cùng một UnifiedCustomer.

\subsection{Thiết kế cơ sở dữ liệu}
\label{subsection:4.2.6}

Cơ sở dữ liệu được thiết kế theo mô hình shared database, trong đó các bảng nghiệp vụ chính có trường \texttt{tenant\_id} để kiểm soát phạm vi dữ liệu. Bảng~\ref{tab:ch4_database_main_tables} tổng hợp các nhóm bảng dữ liệu chính của hệ thống.

\begin{table}[H]
\centering
\caption{Các bảng dữ liệu chính của hệ thống}
\label{tab:ch4_database_main_tables}
\resizebox{\textwidth}{!}{
\begin{tabular}{|p{4cm}|p{5cm}|p{6.5cm}|}
\hline
\textbf{Nhóm dữ liệu} & \textbf{Bảng/khái niệm chính} & \textbf{Vai trò} \\
\hline
Tenant và phân quyền & \texttt{tenants}, người dùng, vai trò/thành viên tenant & Quản lý cửa hàng, tài khoản, quyền truy cập và tenant context. \\
\hline
Product Dataset & \texttt{product\_datasets} & Lưu metadata của dataset sản phẩm sau khi crawl/import và materialize. \\
\hline
KB Artifact & \texttt{product\_dataset\_artifacts} & Lưu artifact đã build từ ProductDataset, gồm artifact path, status, count và quality status. \\
\hline
Tenant-KB Binding & \texttt{tenant\_kb\_bindings} & Liên kết tenant với artifact, hỗ trợ bind/unbind và chính sách cập nhật. \\
\hline
KB Version & \texttt{tenant\_kb\_versions} & Lưu phiên bản KB active hoặc đã publish của tenant, trỏ tới \texttt{kb\_dir}. \\
\hline
Hội thoại & \texttt{conversations}, \texttt{messages} & Lưu conversation, message, chatbot, channel, user external id và trạng thái hội thoại. \\
\hline
Nhận diện khách hàng & \texttt{unified\_customers}, \texttt{customer\_identities}, \texttt{conversations.unified\_customer\_id} & Liên kết các external identity về khách hàng hợp nhất trong phạm vi tenant. \\
\hline
CRM và bán hàng & \texttt{leads}, \texttt{purchase\_requests} & Lưu dữ liệu sau tư vấn, trạng thái xử lý và yêu cầu mua hàng của khách hàng. \\
\hline
General/market layer & \texttt{general\_sources}, \texttt{general\_products}, \texttt{market\_price\_observations} nếu được triển khai đầy đủ & Phục vụ general compare và market price ở phạm vi nâng cấp tiếp theo hoặc module riêng. \\
\hline
\end{tabular}
}
\end{table}

Trong thiết kế này, \texttt{activeKbVersionId} trên tenant là thông tin quan trọng để runtime xác định KB đang được sử dụng. Các bảng binding và version giúp hệ thống tránh phụ thuộc hoàn toàn vào một đường dẫn \texttt{kbDir} cố định. Đồng thời, \texttt{conversations.unified\_customer\_id} giúp kết nối dữ liệu hội thoại với CRM activity theo khách hàng hợp nhất.

\subsection{Thiết kế kiểm soát chất lượng dữ liệu, prompt và kết quả}
\label{subsection:4.2.7}

Dữ liệu đầu vào của RAG cần được kiểm soát trước khi trở thành KB Artifact. Bảng~\ref{tab:ch4_quality_gate_checks} trình bày các kiểm tra chính trong quality gate.

\begin{table}[H]
\centering
\caption{Các kiểm tra chất lượng dữ liệu trước khi build KB Artifact}
\label{tab:ch4_quality_gate_checks}
\resizebox{\textwidth}{!}{
\begin{tabular}{|p{4cm}|p{5cm}|p{6cm}|}
\hline
\textbf{Kiểm tra} & \textbf{Điều kiện kiểm soát} & \textbf{Ý nghĩa} \\
\hline
Manifest count & Số lượng trong manifest phải khớp số dòng dữ liệu thực tế. & Tránh build artifact từ dataset thiếu hoặc lệch số lượng. \\
\hline
Catalog/RAG line count & Số dòng catalog và rag products/chunks phải khớp kỳ vọng. & Bảo đảm dữ liệu truy xuất tương ứng với dữ liệu sản phẩm. \\
\hline
Mojibake trong title/rag title & Nếu title hoặc rag title lỗi mã hóa thì fail. & Tránh đưa tiêu đề lỗi vào prompt và câu trả lời. \\
\hline
Mojibake trong rag text & Có thể warn nếu tỷ lệ nhỏ. & Cho phép demo nếu lỗi nhỏ nhưng vẫn ghi nhận để xử lý tiếp. \\
\hline
Duplicate URL & Tỷ lệ URL trùng vượt ngưỡng thì fail. & Tránh cùng sản phẩm bị nhân bản quá nhiều trong KB. \\
\hline
Price coverage & Độ phủ giá quá thấp thì fail hoặc warn tùy ngưỡng. & Giúp chatbot có đủ dữ liệu khi tư vấn hoặc tham chiếu giá. \\
\hline
Material/dimensions coverage & Độ phủ thấp thì warn. & Ghi nhận hạn chế khi so sánh sản phẩm theo chất liệu/kích thước. \\
\hline
Stale crawl report & Báo cáo crawl quá cũ thì warn hoặc fail theo mức độ. & Tránh dùng dataset không còn phản ánh trạng thái nguồn hiện tại. \\
\hline
Suspicious category & Category nghi ngờ vượt ngưỡng thì warn. & Phát hiện taxonomy chưa ổn định hoặc source cần profile riêng. \\
\hline
\end{tabular}
}
\end{table}

Sau khi dữ liệu đạt điều kiện build, pipeline online vẫn cần kiểm soát prompt, context, token và output. Hệ thống không đưa toàn bộ KB vào prompt. Thay vào đó, chatbot runtime chỉ truy xuất top-k chunk hoặc sản phẩm liên quan nhất với câu hỏi hiện tại, sau đó build prompt theo cấu trúc tách biệt.

Bảng~\ref{tab:ch4_pipeline_control_points} mô tả các điểm kiểm soát trong pipeline RAG và tác động nghiệp vụ.

\begin{table}[H]
\centering
\caption{Các điểm kiểm soát trong pipeline RAG và tác động nghiệp vụ}
\label{tab:ch4_pipeline_control_points}
\resizebox{\textwidth}{!}{
\begin{tabular}{|p{3cm}|p{4cm}|p{4cm}|p{5cm}|}
\hline
\textbf{Bước} & \textbf{Dữ liệu đầu vào} & \textbf{Rủi ro cần kiểm soát} & \textbf{Cơ chế kiểm soát} \\
\hline
Resolve tenant & Token, API request, channel binding & Truy cập sai tenant. & Backend xác định tenant context, tenant admin không được tự chọn tenant khác. \\
\hline
Resolve KB & Tenant, activeKbVersionId & Runtime dùng KB cũ hoặc sai KB. & Ưu tiên ACTIVE\_VERSION, fallback legacy kbDir nếu chưa có active version. \\
\hline
Retrieve context & KB của tenant hoặc corpus theo mode & Context quá nhiều hoặc lẫn dữ liệu. & Truy xuất top-k, giới hạn số chunk/sản phẩm, lọc theo tenant/mode/source. \\
\hline
Build prompt & System rule, mode rule, history, context, user input & Prompt dài, lẫn chỉ dẫn và dữ liệu. & Tách system instruction, mode rule, tenant context, chat history, retrieved evidence, user input và output format. \\
\hline
Generate answer & Prompt đã build & LLM trả rỗng hoặc trả lời ngoài context. & Claude API (Anthropic) kết hợp rule/template fallback và render product answer khi có retrieval context. \\
\hline
Output guard & Câu trả lời và tín hiệu nghiệp vụ & Bịa thông tin sản phẩm/giá hoặc sai format. & Yêu cầu nêu thiếu dữ liệu khi thiếu context, không tự tạo sản phẩm/giá ngoài dữ liệu. \\
\hline
Business action & Mode, conversation, contact, product intent & Tạo lead/purchase sai mode hoặc tạo trùng. & Backend chỉ cho \texttt{tenant\_sales} tạo lead/purchase khi đủ điều kiện và chưa trùng. \\
\hline
\end{tabular}
}
\end{table}

Prompt được tổ chức theo template gồm các phần chính: system instruction, mode rule, tenant context, lịch sử hội thoại cần thiết, retrieved evidence, câu hỏi hiện tại và yêu cầu định dạng đầu ra. Cách tổ chức này giúp mô hình phân biệt rõ đâu là quy tắc, đâu là dữ liệu tham khảo và đâu là yêu cầu của người dùng. Với mode \texttt{market\_price}, hệ thống không để LLM tự bịa khoảng giá; câu trả lời cần dựa trên dữ liệu giá hoặc provider có cấu trúc. Với mode \texttt{general\_compare}, hệ thống không tạo lead hoặc purchase request. Với mode \texttt{tenant\_sales}, backend vẫn kiểm tra điều kiện nghiệp vụ trước khi ghi dữ liệu.

\section{Xây dựng ứng dụng}
\label{section:4.3}

\subsection{Thư viện và công cụ sử dụng}
\label{subsection:4.3.1}

Bảng~\ref{tab:ch4_tools_libraries} trình bày các công cụ và thư viện chính được sử dụng để xây dựng hệ thống.

\begin{table}[H]
\centering
\caption{Danh sách thư viện và công cụ sử dụng}
\label{tab:ch4_tools_libraries}
\resizebox{\textwidth}{!}{
\begin{tabular}{|p{4cm}|p{4cm}|p{7cm}|}
\hline
\textbf{Mục đích} & \textbf{Công cụ/thư viện} & \textbf{Vai trò trong hệ thống} \\
\hline
Backend nghiệp vụ & Java, Spring Boot & Xây dựng API, tenant context, RBAC, Product Dataset, KB binding, conversation, CRM, lead và purchase request. \\
\hline
AI/RAG service & Python, FastAPI & Xử lý chat request, retrieval top-k, prompt building, gọi Claude API (Anthropic) và fallback. \\
\hline
Cơ sở dữ liệu & PostgreSQL & Lưu tenant, conversation, message, Product Dataset, artifact, KB version, customer identity, lead và purchase request. \\
\hline
Frontend & HTML, CSS, JavaScript thuần & Xây dựng admin UI, tenant UI và chat UI dưới dạng static resources được Spring Boot phục vụ trực tiếp. \\
\hline
LLM provider & Claude API & Sinh câu trả lời từ prompt đã được backend và AI service kiểm soát. \\
\hline
Triển khai & Docker Compose & Chạy backend, chatbot service và PostgreSQL trong môi trường demo thống nhất. \\
\hline
Kiểm thử backend & JUnit, Spring Boot Test & Kiểm thử service, controller, tenant isolation, KB version, runtime status và identity. \\
\hline
Kiểm thử AI service & Pytest hoặc script kiểm thử Python & Kiểm thử sales flow, retrieval, fallback và state reset của chatbot runtime. \\
\hline
\end{tabular}
}
\end{table}

\subsection{Kết quả đạt được}
\label{subsection:4.3.2}

Kết quả xây dựng của hệ thống gồm các thành phần chính sau. Thứ nhất là Spring Boot backend, đóng vai trò system-of-record và cung cấp các API nghiệp vụ cho quản trị tenant, chatbot, Product Dataset, KB Artifact, TenantKbBinding, TenantKbVersion, conversation, customer identity, CRM, lead và purchase request. Thứ hai là FastAPI chatbot runtime, đảm nhiệm phần RAG, prompt building, gọi Claude API (Anthropic) và fallback. Thứ ba là PostgreSQL database, lưu dữ liệu nghiệp vụ dùng chung nhưng được kiểm soát theo tenant. Thứ tư là frontend static HTML/CSS/JavaScript, phục vụ các màn hình quản trị và tương tác. Thứ năm là bộ dữ liệu/KB artifact được build từ Product Dataset và bind vào tenant để phục vụ tenant sales.

Về dữ liệu, hệ thống đã hỗ trợ pipeline Product Dataset và KB Artifact thay vì chỉ sử dụng một thư mục KB cố định. Dataset được materialize thành các file phục vụ catalog và RAG, sau đó đi qua quality gate trước khi build artifact. Artifact có thể được bind vào tenant để tạo TenantKbVersion active. Runtime chatbot resolve KB theo active version và có thể được quan sát qua runtime-status.

Về hội thoại, hệ thống hỗ trợ tenant sales, general compare và market price theo các mode khác nhau. Mode \texttt{tenant\_sales} được gắn với tenant-bound KB và có thể tạo lead/purchase request nếu đủ điều kiện. Hai mode \texttt{general\_compare} và \texttt{market\_price} được thiết kế để phục vụ tư vấn chung hoặc tham khảo giá và không tạo tác động nghiệp vụ bán hàng.

Về kênh giao tiếp, hệ thống hỗ trợ Web, Messenger và Telegram. Các kênh này được chuẩn hóa về luồng conversation nội bộ. Cơ chế CustomerIdentityService giúp liên kết các external identity về UnifiedCustomer khi có phone/email đã chuẩn hóa, đồng thời không merge chỉ vì trùng displayName.

\subsection{Minh họa các chức năng chính}
\label{subsection:4.3.3}

Các chức năng chính nên được minh họa bằng ảnh chụp giao diện hoặc ảnh chạy demo trong phụ lục hoặc trong mục này. Nhóm màn hình nên chọn gồm: màn hình quản lý Product Dataset và nguồn KB, màn hình rebuild/publish KB, màn hình runtime-status, màn hình chat tenant sales, màn hình CRM/customer activity và màn hình purchase request.

Khi đưa ảnh minh họa vào bản cuối, mỗi ảnh cần có chú thích ngắn gọn. Ví dụ, màn hình runtime-status nên giải thích các trường \texttt{desired}, \texttt{running}, \texttt{in\_sync}, \texttt{source}, \texttt{kb\_dir} và \texttt{lastUsedAt}. Màn hình chat tenant sales nên thể hiện chatbot đang sử dụng dữ liệu từ active KB version. Màn hình customer activity nên thể hiện các conversation, lead hoặc purchase request được gom theo UnifiedCustomer.

\section{Kiểm thử}
\label{section:4.4}

Kiểm thử của hệ thống tập trung vào các luồng có rủi ro cao trong scope multi-tenant RAG: tách dữ liệu tenant, build và bind KB, runtime active version, RAG tenant sales, cross-channel identity, mode guard và quality gate. Bảng~\ref{tab:ch4_test_evidence_summary} tổng hợp các nhóm kiểm thử và evidence cần đưa vào bản cuối.

\begin{table}[H]
\centering
\caption{Tổng hợp kiểm thử và evidence triển khai}
\label{tab:ch4_test_evidence_summary}
\resizebox{\textwidth}{!}{
\begin{tabular}{|p{2cm}|p{4cm}|p{5cm}|p{5cm}|}
\hline
\textbf{Mã} & \textbf{Nhóm chức năng} & \textbf{Mục tiêu kiểm thử} & \textbf{Evidence/kết quả cần ghi} \\
\hline
TC-01 & Product Dataset/KB Artifact & Kiểm tra dataset được materialize, quality gate chạy và artifact được build. & \texttt{quality\_audit.json}, manifest count, products/chunks count, artifact READY. \\
\hline
TC-02 & Bind KB với tenant & Kiểm tra artifact được bind vào tenant và tạo active KB version. & TenantKbBinding active, TenantKbVersion created, \texttt{activeKbVersionId} được cập nhật. \\
\hline
TC-03 & Runtime ACTIVE\_VERSION & Kiểm tra runtime resolve KB theo active version. & Runtime-status có desired/running cùng source ACTIVE\_VERSION, \texttt{in\_sync=true}. \\
\hline
TC-04 & Runtime evict sau publish/unbind & Kiểm tra publish/unbind làm runtime cũ hết hiệu lực. & Sau publish/unbind: running rỗng hoặc không sync; chat sau đó spawn runtime mới đúng version. \\
\hline
TC-05 & Online tenant\_sales RAG & Kiểm tra chat dùng đúng KB tenant và trả lời dựa trên retrieval context. & Chat smoke với các nhóm sản phẩm như tủ, bàn làm việc, sofa; response có dữ liệu sản phẩm phù hợp. \\
\hline
TC-06 & Cross-channel identity & Kiểm tra Messenger/Telegram/Web có thể liên kết về cùng UnifiedCustomer khi có phone/email. & CustomerIdentityIntegrationTest; same phone/email tạo cùng unifiedCustomerId. \\
\hline
TC-07 & Không merge displayName-only & Kiểm tra hệ thống không gộp khách chỉ vì trùng tên hiển thị. & Test displayName-only không merge. \\
\hline
TC-08 & Tenant isolation trong identity & Kiểm tra cùng phone/email nhưng khác tenant không bị trộn dữ liệu. & Test different tenant không merge. \\
\hline
TC-09 & CRM activity & Kiểm tra truy vấn lịch sử theo unifiedCustomerId. & CrmCustomerActivityServiceTest; API activity trả conversation/lead/purchase đúng tenant. \\
\hline
TC-10 & Mode guard & Kiểm tra general\_compare và market\_price không tạo purchase request. & Test mode control; chỉ tenant\_sales mới được tạo purchase khi đủ điều kiện. \\
\hline
TC-11 & Prompt/context/token control & Kiểm tra chatbot không đưa toàn bộ KB vào prompt và dùng top-k context. & Log hoặc test retrieval top-k, prompt template, fallback khi provider rỗng. \\
\hline
TC-12 & Reset/new session & Kiểm tra reset/new session không làm mất dữ liệu đã hoàn tất nhưng xóa trạng thái đang xử lý. & Test reset conversation, clear Python state, discard in-progress lead/purchase. \\
\hline
\end{tabular}
}
\end{table}

Trong bản cuối, các dòng ở Bảng~\ref{tab:ch4_test_evidence_summary} nên được điền bằng tên test, script hoặc file evidence cụ thể. Các evidence phù hợp gồm \texttt{tmp/verify\_admin\_kb\_flow.py}, runtime-status sau publish, runtime-status sau chat, \texttt{kb\_lazy\_verify\_summary.json}, \texttt{CustomerIdentityIntegrationTest}, \texttt{CrmCustomerActivityServiceTest}, \texttt{LlmInstanceManagerTest}, \texttt{TenantKbVersionServiceTest}, \texttt{TenantKbRuntimeStatusControllerTest} và các file quality audit của dataset.

\section{Triển khai}
\label{section:4.5}

Hệ thống được triển khai thử nghiệm theo mô hình nhiều service. Spring Boot backend chạy API nghiệp vụ và phục vụ frontend static. FastAPI chatbot service chạy riêng để xử lý RAG và gọi Claude API (Anthropic) hoặc dùng rule fallback. PostgreSQL được sử dụng làm cơ sở dữ liệu chung. Các thành phần này có thể được khởi chạy bằng Docker Compose trong môi trường demo.

Trong mô hình triển khai, backend cần cấu hình kết nối tới PostgreSQL, đường dẫn tới chatbot service, API key cho Claude API (Anthropic) và đường dẫn lưu KB artifact. Chatbot service cần có quyền đọc thư mục KB tương ứng với active KB version của tenant. Khi backend publish hoặc unbind KB, runtime của tenant bị evict; request chat tiếp theo sẽ làm runtime load lại KB mới.

Runtime-status được sử dụng như một cơ chế quan sát khi triển khai. Các trường \texttt{desired}, \texttt{running}, \texttt{process\_alive}, \texttt{source}, \texttt{kb\_dir}, \texttt{version\_tag}, \texttt{in\_sync} và \texttt{lastUsedAt} cho biết runtime đang chạy có khớp với KB version mong muốn hay không. Đây là thông tin quan trọng để kiểm tra sau khi rebuild, publish, bind hoặc unbind KB.

Với frontend, các trang admin và tenant được đóng gói trong thư mục static resources của Spring Boot, nên không cần một frontend build service riêng. Cách triển khai này làm giảm số thành phần phải vận hành trong phạm vi đồ án, đồng thời vẫn đủ để demo các chức năng chính như quản lý dataset, publish KB, chat, runtime-status, CRM activity và purchase request.

\end{document}